% ============================================================
\section{Phần IV: Đánh giá hệ thống}
\label{sec:eval}
% ============================================================

\subsection{Tập kiểm thử}

Nhóm đánh giá trên ba bộ ở \cref{tab:testsets}. Bộ thử thách nhắm vào những trường hợp bộ
của đề còn ít: cách diễn đạt khác, câu KB không có dữ kiện, thực thể không tồn tại và câu
ngoài phạm vi.

\begin{table}[htbp]
  \centering
  \small
  \begin{tabularx}{\linewidth}{|P{3.4cm}|c|L|}
    \hline
    \textbf{Bộ} & \textbf{Số câu} & \textbf{Thành phần} \\ \hline
    Bộ của đề & 42 & 6 thông tin môn học, 4 lịch học, 11 chủ đề, 2 LO, 6 bài tập lớn,
      1 deadline, 3 tài liệu, 3 quy định, 3 câu ngữ cảnh, 3 câu ngoài phạm vi \\ \hline
    Bộ thử thách & 22 & 6 cách diễn đạt khác, 3 câu KB không có dữ kiện, 6 alias và tính
      năng BTL, 4 thực thể không tồn tại, 3 câu ngoài phạm vi \\ \hline
    Hội thoại của đề & 6 phiên & Câu nêu thực thể, rồi câu tham chiếu tới thực thể đó \\ \hline
  \end{tabularx}
  \caption{Các bộ kiểm thử}
  \label{tab:testsets}
\end{table}

\subsection{Phương pháp chấm}

Một câu đúng end-to-end khi intent, entity, trạng thái truy vấn đều đúng và câu trả lời chứa
đủ các fact bắt buộc. Câu ngữ cảnh được chấm intent trước khi giải và entity sau khi giải;
câu ngoài phạm vi chỉ đúng khi parser từ chối và hệ thống không truy vấn KB. Câu trả lời
được chấm theo fact thay vì so nguyên văn để không phụ thuộc template. Ngữ cảnh được xóa
trước mỗi bộ và giữa các hội thoại.

\subsection{Kết quả}

\begin{table}[htbp]
  \centering
  \small
  \begin{tabular}{|l|r|r|r|r|r|r|}
    \hline
    \textbf{Bộ} & \textbf{Số câu} & \textbf{Intent} & \textbf{Entity} &
      \textbf{Query} & \textbf{Answer} & \textbf{End-to-end} \\ \hline
    Bộ của đề & 42 & 42 & 42 & 42 & 42 & 42 (100\%) \\ \hline
    Bộ thử thách & 22 & 20 & 20 & 20 & 20 & 20 (90,91\%) \\ \hline
    \textbf{Tổng} & \textbf{64} & 62 & 62 & 62 & 62 & \textbf{62 (96,88\%)} \\ \hline
    Hội thoại & 6 phiên & \multicolumn{5}{c|}{6/6 phiên giải đúng thực thể tham chiếu} \\ \hline
  \end{tabular}
  \caption{Kết quả đánh giá}
  \label{tab:results}
\end{table}

Hai câu sai đều sai ngay từ tầng cú pháp; các tầng sau không thêm lỗi. Trong 7 câu mong đợi
``không tìm thấy'', 6 câu đạt, câu còn lại là C007.

\subsection{Phân tích các trường hợp sai}

\begin{table}[htbp]
  \centering
  \small
  \begin{tabularx}{\linewidth}{|P{3.5cm}|P{3cm}|L|L|}
    \hline
    \textbf{Câu} & \textbf{Mong đợi} & \textbf{Nguyên nhân} & \textbf{Hướng sửa} \\ \hline
    C007: ``Môn NLP học mấy buổi mỗi tuần?'' & Hiểu câu, báo KB không có dữ kiện &
      Văn phạm chưa có trường số buổi học nên câu bị từ chối thay vì báo thiếu dữ kiện. &
      Thêm trường này vào nhánh thông tin môn học. \\ \hline
    C008: ``Tuần thứ ba tụi mình học bài gì vậy?'' & Lịch học tuần 3 &
      Văn phạm chưa nhận số thứ tự viết bằng chữ như ``thứ ba'', và chưa có cụm ``tụi
      mình'', ``bài gì''. & Thêm số thứ tự vào văn phạm và từ điển thực thể. \\ \hline
  \end{tabularx}
  \caption{Các câu trả lời sai}
  \label{tab:errors}
\end{table}

Cả hai lỗi cùng một gốc: văn phạm đóng không nhận cách diễn đạt chưa được viết thành luật,
dù nghĩa của câu nằm trong miền. Đây là đánh đổi của hướng Classical: mất độ phủ để đổi lấy
khả năng giải thích. Nhóm giữ nguyên hai ca này thay vì thêm luật riêng cho đúng hai câu
kiểm thử.

Các câu hỏi tuần 99, chương 13, LO9.9 được trả lời ``không tìm thấy'' là đúng vì KB không có
các thực thể đó. Câu có token thừa hoặc hai câu hỏi trong một lượt bị từ chối theo thiết kế.

\subsection{Kiểm thử tự động}
\label{sec:testing}

Ngoài đánh giá end-to-end, mỗi tầng được kiểm thử riêng trên đầu vào chuẩn sinh từ bản cài
đặt tham khảo, nên lỗi được khoanh vùng đúng tầng. Bộ test còn có các ca hồi quy như luật
rỗng, đệ quy trái, mơ hồ, token thừa và thiếu dữ kiện KB; tổng cộng 49 test, đều pass. Cùng
với việc kiểm tra văn phạm ở \cref{sec:grammarcheck}, mọi thay đổi đều được kiểm lại trên
toàn bộ câu kiểm thử và 1.000 câu sinh.
